% Supplement, part a: model notes S1-S4.
% Source paths in "% src:" comments.
%   src/..., tests/... and Cargo.toml are relative to /home/guarracino/git/impg at HEAD 0987752 (impg 0.5.0).
%   T: is short for tests/test_transitive_integrity.rs.
%   sweepga@4713ce0 is /home/guarracino/.cargo/git/checkouts/sweepga-d4f4c2d347bec45b/4713ce0 (Cargo.toml:71).
%   seqwish@0a7e580 is /home/guarracino/.cargo/git/checkouts/seqwish-be7416204a0fe48d/0a7e580 (Cargo.lock:2872-2874).
%   seqwish-cpp is /home/guarracino/git/seqwish/cpp (checkout 97d603f).
%   "impg <cmd> --help" refers to /tmp/claude-1000/-home-guarracino-git-impg-paper/ed3cfa77-35b3-454f-91de-3d0a63c6ad53/scratchpad/phase2/plan_help_<cmd>.txt.
%   VERIFY_CORRECTIONS.md is /tmp/claude-1000/-home-guarracino-git-impg-paper/ed3cfa77-35b3-454f-91de-3d0a63c6ad53/scratchpad/phase1/VERIFY_CORRECTIONS.md.
% Nothing in this note was run. "Asserted" means a test checks the value, and "derived" means we followed the code by hand.
% In S2, S3 and S4 the line numbers in parentheses are the sources of each step.

\subsection*{Supplementary Note S1: worked example}

We trace fixtures of tests/test\_transitive\_integrity.rs through Section \ref{sec:model}, with every input coordinate copied from the test.
% src: T:1-915
A value marked asserted is checked by the test, and a value marked derived follows from the code and was not run.

\textbf{Two entries per record.}
The record \texttt{A 1000 0 100 + B 1000 200 300 \dots{} cg:Z:100=} enters the tree of B and, as a REVERSED entry, the tree of A.
% src: T:231; src/impg.rs:1666-1709
A:0-100 projects to B:200-300, and B:200-300 projects to A:0-100 (both asserted).
% src: T:237-262, 264-289

\textbf{An insertion.}
The record \texttt{A 1000 0 110 + B 1000 0 100 \dots{} cg:Z:50=10I50=} leaves A:50-60 without a partner in B.
% src: T:457-462
The REVERSED entry in the tree of A reads the inverted CIGAR \texttt{50=10D50=}.
% src: src/impg.rs:142-162, 651-654; unit tests src/impg.rs:3232-3290
A:0-50 projects to B:0-50 and A:60-110 to B:50-100 (derived).
% src: hand trace of src/impg.rs:2791-2929
The test asserts that the B ends next to the insertion lie in 45 to 55.
At the outer ends, it asserts only a start of at most 5 and an end of at least 95.
% src: T:492-497 (b_start <= 5, 45 <= b_end <= 55), 523-528 (45 <= b_start <= 55, b_end >= 95)
A:45-65 projects to B:45-55, so 20 bases of A give 10 bases of B (derived).
% src: src/impg.rs:2853-2871 (deletion columns add no query base)
In the other direction, B:0-100 projects to A:0-110, because the hull includes the \texttt{10I} run (derived).
% src: src/impg.rs:2838-2852

\textbf{Reverse strand.}
For \texttt{A 1000 0 100 - B 1000 0 100 \dots{} cg:Z:100=}, the walk for A:0-50 starts at B position 100 and moves down.
% src: T:302; src/impg.rs:2808-2813, 2873-2898
It returns the pair (100, 50), which the BED writer prints as B:50-100 on strand $-$ (derived).
% src: src/main.rs:12074-12078
The test asserts only that the midpoint of the B row is at least 50.
% src: T:333-341

\textbf{A two-hop closure.}
Two \texttt{cg:Z:100=} records, \texttt{A 1000 0 100 + B 1000 0 100} and \texttt{B 1000 0 100 + C 1000 0 100}, form a chain.
% src: T:160-163
The test queries A:25-75 at $m=2$, $\ell=0$ and $\delta=10$.
% src: T:170-186; src/main.rs:4236, 4247
At depth 0 the tree of A holds one entry, the REVERSED copy of the first record, which gives B:25-75 for the next frontier.
% src: src/impg.rs:2405-2441, 1686-1712, 2574-2588
At depth 1 the tree of B gives A:25-75, already visited, and C:25-75, which is new.
% src: src/impg.rs:1666-1709, 2578-2588
The loop then stops at $k=m=2$, and BED output with \texttt{-d 0} merges the two A rows (derived).
% src: src/impg.rs:2444; src/main.rs:12055-12071
The test asserts 45 to 55 bp per row, with B and C starting at 20 to 30 and ending at 70 to 80.
% src: T:192-217
On the chain A, B, C, D it asserts that $m=1$ returns A and B only, and that $m=2$ adds C but not D.
% src: T:693-698, 728-735, 761-766

\textbf{Snapping adds reach that no alignment supports.}
The 5 bp gap fixture aligns C:0-1000 and B:0-1000 to A:0-1000, C:1005-2005 to B:0-1000, D:0-300 to C:1605-1905 and E:0-4 to C:1000-1004.
% src: T:786, 795-800
The test queries A:0-1000 at $m=0$, $\ell=101$ and $\delta=10$.
% src: T:809-816; src/main.rs:4247, 4253-4255
At depth 1, B:0-1000 projects to C:1005-2005, which starts 5 bp after the visited end 1000.
% src: T:798; src/impg.rs:2452-2533 (derived)
Since 5 is less than 10, the new part snaps to C:1000-2005 and enters the frontier.
% src: src/impg.rs:296-298, 2578-2588 (derived)
At depth 2 this part overlaps the E record on C:1000-1004, so E:0-4 is reported, and D:0-300 is reached as well.
% src: src/impg.rs:2466-2471, 2559-2571 (derived)
The test asserts D and E:0-4, and it asserts that the C rows stay exactly C:0-1000 and C:1005-2005.
% src: T:820-844
With $\delta=0$, or with a 15 bp gap at $\delta=10$, it asserts that E is not reported.
% src: T:787, 825-830, 846-858
No projection covers C:1000-1005 and no alignment chain links E to A, so snapping alone put E into the closure.
% src: T:795-800

\textbf{Strand over two reverse hops.}
No fixture has two reverse-strand hops, so this case is derived only.
% src: T:298-342 (one reverse hop only)
With both records of the chain on strand $-$, A:25-75 projects to the pair (75, 25) on B, and $V_B$ keeps B:25-75 without orientation.
% src: src/impg.rs:282-286, 374-379, 2808-2813
B:25-75 then gives the pair (75, 25) on C, so the BED row reads C:25-75 on strand $-$.
After two reverse complements, C:25-75 has strand $+$ relative to A.
% src: src/main.rs:12074-12078
\authornote{Decision D3: this case is the fixture that code fix C7 needs, whichever semantics the authors choose.}

\subsection*{Supplementary Note S2: proofs}

$V_x$ is the visited set of sequence $x$ when Algorithm \ref{alg:bfs} stops, $V$ is the union of the $V_x$, and line numbers refer to src/impg.rs.

\textbf{Assumptions of Proposition \ref{prop:closed}.}
(A1) One combined index, built with bidirectional insertion from every record of $A$ (lines 1686-1712), and not the per-file mode (src/multi\_impg.rs:813-952).
(A2) The default breadth-first search, not \texttt{-{}-transitive-dfs} (src/main.rs:4229-4232).
(A3) Exact projection of PAF records with CIGARs.
For TPA and 1ALN records, $a$ is the alignment of the rebuilt CIGAR, and no hit is dropped.
A segment also rebuilds to the same columns whatever range triggers it.
\authornote{We did not check this last part of (A3) in the code. E1 tests its consequence on the release data.}
(A4) $m=0$, $\ell\le 1$ and $\delta=0$.
(A5) No identity filter, subset filter, mask or \texttt{-l} (lines 2553-2557).
(A6) Sequences shorter than $2^{31}$ bp and CIGAR operations shorter than $2^{29}$ bp (lines 81-109, 164-174).
(A7) $|I|\ge\ell$, which the CLI enforces (src/main.rs:7475-7480, 10598-10614).

\textbf{Lemma 1 (hull).}
Let $r$ be an entry on sequence $x$ and $J$ a non-empty interval inside the span of $r$ on $x$.
If an aligned column of $r$ joins $b\in J$ to $c$, then $\pi_r(J)$ is not empty and contains $c$.

\emph{Proof.}
The walk visits the operations in order and stops only after passing the end of $J$ (lines 2829-2835).
The operation that holds the column of $b$ overlaps $J$, so it extends the hull over its overlap, which includes $c$ (lines 2873-2894).
Query positions move in one direction along the walk, so the hull from the first to the last recorded position contains $c$.
That operation adds at least one base on each side, so the tests at lines 2905-2907 pass.
A REVERSED entry first inverts the CIGAR (lines 142-162, 651-654), which keeps the same columns with query and target swapped (unit tests at lines 3232-3290).
$\square$

\textbf{Lemma 2 (frontier coverage).}
Under (A4), (A5) and (A7), every base of $V$ lies in a frontier interval that the loop processes.

\emph{Proof.}
Without masks, $I$ enters $V_j$ and the first frontier whole (lines 2399-2441).
A later projection $K$ onto $y\neq x$ enters $V_y$ through \texttt{insert}, because $\delta=0$ (lines 276-278).
With distance 0 in every visited set (lines 2145-2159), \texttt{insert} snaps to nothing.
It returns exactly the parts of $K$ not yet in $V_y$ (lines 374-457).
Each part has at least one base, so $\ell\le 1$ puts it in the next frontier (line 2585), and merging keeps the union (lines 2597-2615).
With $m=0$ the loop runs until the frontier is empty (line 2444).
$\square$

\textbf{Termination.}
Each non-empty frontier holds only bases new to $V$, and $V$ never shrinks, so the loop runs at most $1+\sum_i|s_i|$ depths.

\textbf{Proof of Proposition \ref{prop:closed}.}
Let $b\in V_x$, and let $r\in A$ between $x$ and $y\neq x$ align $b$ to $c$.
By Lemma 2, $b$ lies in a processed frontier interval $(x,J)$.
By (A1), the tree of $x$ holds a copy of $r$, forward if $x$ is its target and REVERSED if $x$ is its query.
The tree query keeps every entry with a non-empty overlap with $J$ (lines 2460-2471), and this overlap contains $b$.
By Lemma 1 the projection contains $c$, and by (A5) no filter removes it.
Since $y\neq x$, line 2575 inserts it into $V_y$, so $c\in V$.
$\square$

\textbf{Corollary.}
If no record of $A$ joins a sequence to itself, every match joins two different sequences.
Induction on chain length then shows that $V$ contains the $\equiv_A$ class of every base of $I$.
Every reported row is then $I$ or a projection that line 2575 inserted with $\delta=0$, so $H_0(I)=V$.
The containment is strict in general, because projections include insertions and mismatches.

\textbf{What fails without each assumption.}
With $m>0$, bases more than $m$ hops away are missing.
With $\ell>1$, a part shorter than $\ell$ is visited but never queried (line 2585).
With $\delta>0$, snapping adds fewer than $\delta$ uncovered bases at each end of a projection and can drop fewer than $\delta$ bases before the end of a part (lines 296-303, 321-335).
A filtered hit is neither written nor expanded, and masked bases never enter a frontier.
With \texttt{-l}, a row shorter than the limit still enters $V$ but is not written, so $H_0(I)$ can miss bases of $V$ (lines 2553-2588).
With \texttt{-{}-unidirectional}, the tree of $x$ misses $r$ when $x$ is the query of $r$.
Self-hits are written but not expanded (line 2575).
With \texttt{-{}-approximate} the boundary is interpolated (lines 1471-1520), and Lemma 1 can fail.
A dropped tracepoint hit removes its columns (lines 1230-1236, 1247-1254, 1330-1338).

\textbf{Assumptions of Proposition \ref{prop:restrict}.}
(B1) A seqwish match is an \texttt{M}, \texttt{=} or \texttt{X} column with equal bases, not \texttt{N}, at different positions (seqwish@0a7e580 src/alignments.rs:133, 202-204, and seqwish-cpp src/alignments.cpp:54-83).
(B2) seqwish runs with \texttt{-k} at most 1 and with \texttt{-r}, \texttt{-l} and \texttt{-f} at their default 0.
Each of these four options can depend on context outside $C$ (seqwish \texttt{-{}-help}).
(B3) seqwish puts two oriented positions in one node class exactly when the closure of its matches makes them equivalent \citep{Garrison2023seqwish}.
\authornote{(B3) is the definition in the seqwish paper. We did not check in the seqwish code that the batched transitive closure (-B) gives the same classes for every batch size.}
(B4) $C$ is closed under $R(A)$.
(B5) $A|_C$ has exactly the aligned columns of $A$ with both bases in $C$, with the same bases and strands.
(B6) Both runs read the same sequences.
We compare node classes, not node identifiers or node boundaries, because compaction can cut nodes differently at the ends of $C$.

\textbf{Proof of Proposition \ref{prop:restrict}.}
Let $Q_C$ be the oriented positions of the bases of $C$.
By (B1) and (B5), $R(A|_C)=R(A)\cap(Q_C\times Q_C)$, and by (B2) seqwish keeps every match in both runs.
For $u,v\in Q_C$, a chain of matches of $A|_C$ is a chain of matches of $A$, so $u\equiv_{A|_C}v$ implies $u\equiv_A v$.
Conversely, by (B4) and induction, every link of a chain of matches of $A$ from $u$ stays in $Q_C$ and so lies in $R(A|_C)$.
Hence the two relations agree on $Q_C$, every class of $\equiv_A$ that meets $Q_C$ lies inside it, and (B3) gives the claim.
$\square$

\textbf{Why \texttt{-k} must be at most 1.}
Closedness allows a run of 5 matches with 2 columns outside $C$ and 3 inside, cut at the same column on both sequences.
At \texttt{-k 4}, seqwish keeps the run of 5 in $G(A)$ and drops the run of 3 in $G(A|_C)$ (seqwish@0a7e580 src/alignments.rs:141).

\textbf{Link between the propositions.}
Under (A1) to (A7) with no self-alignment, $V$ is closed under $R(A)$, so Proposition \ref{prop:restrict} applies with $C=V$.
With \texttt{-{}-no-merge} and \texttt{-o paf}, every processed frontier interval $J$ and entry $r$ give one row with the CIGAR slice of $r$ over $J$.
% src: lines 2509-2523, 2559-2571; src/main.rs:12195-12295
By Lemma 2 and Proposition \ref{prop:closed}, the matches of these rows are exactly $R(A)\cap(Q_V\times Q_V)$, and columns written twice do not change the closure.
seqwish on these rows should therefore give the classes of $G(A)$ over $V$.
We derived this from the code, no test covers PAF output on the alignment backend, and \texttt{-o gfa} does not use these rows.
% src: VERIFY_CORRECTIONS.md (query output formats); src/main.rs:12324-12340
\pending{E6}{bases of V whose node class differs between seqwish on the closure PAF and the global seqwish graph, on yeast chrV and the C4 holdout}

\subsection*{Supplementary Note S3: sparsification strategies as implemented}

\texttt{impg align -{}-sparsify} chooses which pairs of genomes to align, and its default \texttt{none} keeps all pairs.
% src: impg align --help (--sparsify, default none); sweepga@4713ce0 src/cli.rs:313-319
The rules live in sweepga \citep{Garrison2026sweepga} at commit 4713ce0, and the library default of impg, \texttt{giant:0.99}, does not reach the command line.
% src: Cargo.toml:71; src/commands/align.rs:53-57 (AlignConfig::default); src/main.rs:9518 (CLI value)
For PanSN input the unit is the haplotype, and $n$ counts haplotypes.
% src: sweepga@4713ce0 src/knn_graph.rs:636-679
A selected haplotype pair expands to all contig pairs across it, and contig pairs within one haplotype are always kept.
% src: sweepga@4713ce0 src/knn_graph.rs:584-615
Distances are Mash distances \citep{Ondov2016mash} between merged haplotype sketches, with k-mer size 15 and sketch size 1,000 by default.
% src: sweepga@4713ce0 src/knn_graph.rs:561-583, 647-656; impg align --help (--mash-kmer-size 15, --mash-sketch-size 1000)

\begin{itemize}
\item \texttt{none} keeps all $n(n-1)/2$ pairs.
% src: sweepga@4713ce0 src/knn_graph.rs:505-513
\item \texttt{random:}$f$ keeps a pair when a hash of its two indices is at most $f\,(2^{64}-1)$, with Rust's \texttt{DefaultHasher} and no user seed.
% src: sweepga@4713ce0 src/knn_graph.rs:60-66, 371-392
\item \texttt{tree:}$k_n$\texttt{:}$k_f$\texttt{:}$r$ links each unit to its $k_n$ nearest and $k_f$ farthest units and adds random pairs at fraction $r$.
% src: sweepga@4713ce0 src/knn_graph.rs:94-133, 243-282, 343-368
\item \texttt{giant:}$p$ sets $f=\min\{1,\lceil n\ln n\,(-\ln p)/2\rceil/\binom{n}{2}\}$ and $k=\max\{2,\lceil fn\rceil\}$, then applies \texttt{tree:}$k$\texttt{:1:0.01}.
% src: sweepga@4713ce0 src/knn_graph.rs:534-544 (sketch path that impg align uses), 455-472
At $p=0.99$, $k=2$ for $n$ = 7, 235, 466, 1,000 and 10,000.
% src: python3 arithmetic in this step (target edges 1, 7, 15, 35 and 463); VERIFY_CORRECTIONS.md (partition-regional verifier notes)
The rule does not sample an Erd\H{o}s-R\'enyi graph \citep{Erdos1959random}, and no code checks that the pairs connect all units.
% src: sweepga@4713ce0 src/knn_graph.rs:534-544; sweepga@4713ce0 src/lib.rs:2-28 does not declare giant_component.rs, so that module is not compiled
\item \texttt{auto} keeps all pairs for $n\le 10$ and applies \texttt{giant:0.99} for $11\le n\le 50$ and \texttt{tree:5:2:0.05} for $n>50$.
% src: sweepga@4713ce0 src/knn_graph.rs:514-531
\item \texttt{wfmash:}$f$ keeps all pairs and passes a mapping density to wfmash \citep{Guarracino2024wfmash}, with $10\ln(n)/n$ for $n$ PanSN haplotypes under \texttt{wfmash:auto}.
% src: sweepga@4713ce0 src/knn_graph.rs:205-216, 549-557, src/orchestrator.rs:10-21; src/commands/align.rs:421, 478-481; src/commands/mod.rs:102-105 (num_genomes counts haplotypes); src/main.rs:10435-10440 (rejected with other aligners)
\end{itemize}

impg aligns every pair of input files that holds a selected pair, and only wfmash also receives the sequence pairs.
% src: src/commands/align.rs:434-440 (collect_file_pairs), 483-500 (pairs file only for wfmash)
With FastGA, the pair set takes full effect only when each haplotype has its own FASTA file, and no effect with one FASTA file.
% src: src/commands/align.rs:377-396 (collect_file_pairs), 485 (pairs file only for wfmash), 524-530 (one input file is self-aligned whole); sweepga@4713ce0 src/knn_graph.rs:636-657 (haplotype unit); VERIFY_CORRECTIONS.md (partition-regional)
impg tests the parser of these rules and the rejection of \texttt{wfmash:auto} with another aligner, and no impg test runs a sparsified pair set.
% src: src/commands/align.rs:878-898; tests/test_aligner_validation.rs:22-31; sweepga@4713ce0 src/knn_graph.rs:680-1000 (sweepga's own unit tests)

\subsection*{Supplementary Note S4: partition termination}

We consider \texttt{partition} on the alignment backend with breadth-first search, and line numbers refer to src/commands/partition.rs.
Rounds repeat while windows remain (line 297), and each round processes its windows and then selects new ones among unowned bases (lines 297-618).

\textbf{Claim.}
For any window size $w\ge 1$ and depth $m$, the loop ends after at most $1+\sum_i|s_i|$ rounds.

\emph{Proof.}
Both window loops set the window end to $\min(pos+w,\,end)$ and move $pos$ there (lines 229-246, 954-968).
For $w\ge 1$ each step advances $pos$ by at least 1, so each loop ends.
Each closure ends after $m$ depths, or by the termination argument of Note S2 when $m=0$.
Every window joins its own closure output (lines 421, 981-997).
\texttt{mask\_and\_update\_regions} masks that output and removes it from the missing set (lines 1308-1313, 1315-1376).
Windows come only from sequences with missing bases (lines 764-944), and a sequence leaves the missing set once it has none (lines 1373-1375).
So a round with a window removes at least one missing base, and a round without one ends the loop (lines 946-948).
$\square$

\textbf{Window size 0.}
With $w=0$ the window end equals $pos$, and the tail test $0<0$ fails (lines 232-246, 957-967).
The loop appends the empty window $(pos,pos)$, never advances, and grows the window list until memory runs out (derived, not run).
The command accepts \texttt{-w 0}, because the only window-size check in src/main.rs belongs to \texttt{impg infer} (src/main.rs:4771-4774, 8644-8649).
Code fix C11 of the plan adds the check, and values above 2,147,483,647 are cast to 32 bits without a check (lines 232, 957).
% src: plan.md section 9 (C11)
Termination says nothing of whether every base ends in exactly one partition.
\pending{E9}{partitions counted from the BED files, uncovered bp and bp owned by more than one partition on ScRAP 142}
